\section*{第 \ref{ch:g9:trig} 章 --- 直角三角形中的三角学}

\begin{solution}{exo:g9:trig:1}
斜边：$\sqrt{9^2 + 12^2} = \sqrt{81 + 144} = \sqrt{225} = 15$。

另一直角边：$\sqrt{17^2 - 8^2} = \sqrt{289 - 64} = \sqrt{225} = 15$。
\end{solution}

\begin{solution}{exo:g9:trig:2}
$7^2 + 24^2 = 49 + 576 = 625 = 25^2$：是直角三角形（勾股定理的逆定理，
直角所对的边为 $25$）。

$5^2 + 6^2 = 61 \neq 64 = 8^2$：不是直角三角形。
\end{solution}

\begin{solution}{exo:g9:trig:3}
斜边为 $[EF]$（与 $D$ 处的直角相对）。
与 $\theta = \widehat E$ 相对的边为 $[DF]$；邻边为
$[DE]$。由此可得
\[
\cos\theta = \frac{DE}{EF}, \qquad
\sin\theta = \frac{DF}{EF}, \qquad
\tan\theta = \frac{DF}{DE}.
\]
\end{solution}

\begin{solution}{exo:g9:trig:4}
对边：$10 \sin 28^\circ \approx 10 \times 0.469 = 4.7$。
邻边：$10 \cos 28^\circ \approx 10 \times 0.883 = 8.8$。
（验算：$4.7^2 + 8.8^2 \approx 22 + 77 \approx 100$。）
\end{solution}

\begin{solution}{exo:g9:trig:5}
$\tan\theta = \dfrac{4.5}{6} = 0.75$，于是
$\theta = \tan^{-1}(0.75) \approx 37^\circ$。
\end{solution}

\begin{solution}{exo:g9:trig:6}
在由观察者、地面上位于无人机正下方的点与无人机
构成的直角三角形中，高度 $120$ m 为仰角所对的边，
水平距离 $d$ 为邻边。
于是
\[
\tan 32^\circ = \frac{120}{d},
\qquad
d = \frac{120}{\tan 32^\circ} \approx \frac{120}{0.625} = 192 \text{ m}.
\]
\end{solution}

\begin{solution}{exo:g9:trig:7}
升高（$1.2$ m）为角所对的边，坡道长 $L$ 为
斜边：
\[
\sin 5^\circ = \frac{1.2}{L},
\qquad
L = \frac{1.2}{\sin 5^\circ} \approx \frac{1.2}{0.0872} \approx 13.8
\text{ m}.
\]
坡道长度至少为 $13.8$ m。
\end{solution}

\begin{solution}{exo:g9:trig:8}
由 $\cos^2\theta + \sin^2\theta = 1$ 得：
$\sin^2\theta = 1 - 0.36 = 0.64$，且锐角的 $\sin\theta > 0$，
故 $\sin\theta = 0.8$。于是
$\tan\theta = \dfrac{\sin\theta}{\cos\theta} = \dfrac{0.8}{0.6} =
\dfrac43$。
\end{solution}

\begin{solution}{exo:g9:trig:9}
\emph{等腰直角三角形}，两直角边为 $1$，斜边为
$\sqrt{1 + 1} = \sqrt2$；每个锐角为 $45^\circ$，且
\[
\cos 45^\circ = \frac{1}{\sqrt2} = \frac{\sqrt2}{2} = \sin 45^\circ,
\qquad
\tan 45^\circ = \frac11 = 1 .
\]

\emph{等边三角形}，边长为 $1$，其高将它分成两个
直角三角形，斜边为 $1$，底为 $\frac12$，高为
$h = \sqrt{1 - \frac14} = \frac{\sqrt3}{2}$（勾股定理）。其中各角为
$60^\circ$（底边处）与 $30^\circ$（顶点处）。由此读出
各比值：
\[
\cos 60^\circ = \frac{1/2}{1} = \frac12, \quad
\sin 60^\circ = \frac{\sqrt3/2}{1} = \frac{\sqrt3}{2}, \quad
\cos 30^\circ = \frac{\sqrt3}{2}, \quad
\sin 30^\circ = \frac12 .
\]
\end{solution}

\begin{solution}{pb:g9:trig:1}
\textbf{1.} 视线以上的高度是
$35^\circ$ 所对的边，邻边为 $60$ m：
$h = 60 \tan 35^\circ \approx 60 \times 0.700 = 42$ m。

\textbf{2.} 由 $\tan\theta = 0.12$ 得
$\theta = \tan^{-1}(0.12) \approx 6.8^\circ$ --- 对道路来说算陡，
但角度并不大。一条 $45^\circ$ 的道路
满足 $\tan 45^\circ = 1$，即坡度为 $100\,\%$。（“
百分之一百”远非竖直——这是
一个与滑雪缆车相关的常见误解。）

\textbf{3.} 在两锐角分别为 $\theta$ 与
$90^\circ - \theta$ 的直角三角形中，一锐角的对边
恰是另一锐角的邻边，斜边则为公共边。于是
$\sin(90^\circ - \theta) = \frac{\text{其对边}}{h} =
\frac{\text{邻边 }\theta}{h} = \cos\theta$，且
同理，$\cos(90^\circ - \theta) = \sin\theta$：
\emph{余}弦是\emph{余}角的正弦。

\textbf{4.} 高度：$8 \sin 60^\circ = 8 \times
\frac{\sqrt3}{2} = 4\sqrt3 \approx 6.93$ m。底端距离：
恰为 $8 \cos 60^\circ = 4$ m。

\textbf{5.} $0.574^2 + 0.819^2 = 0.329 + 0.671 = 1.000$（按
四舍五入）；且 $\left(\frac12\right)^2 +
\left(\frac{\sqrt3}{2}\right)^2 = \frac14 + \frac34 = 1$。任一
组数值不满足恒等式，说明读数有误，或
计算器模式设置错误 --- 一个免费的错误检测器。

\textbf{6.} $\tan 40^\circ = \dfrac hx$ 与
$\tan 30^\circ = \dfrac{h}{x + 200}$。

\textbf{7.} $x = \dfrac{h}{\tan 40^\circ}$ 与
$x + 200 = \dfrac{h}{\tan 30^\circ}$；两式相减，
$200 = h\left(\frac{1}{\tan 30^\circ} -
\frac{1}{\tan 40^\circ}\right)$，由此得公式。代入数值：
$\frac{1}{0.577} \approx 1.732$，$\frac{1}{0.839} \approx
1.192$，差为 $0.540$：
$h \approx \frac{200}{0.540} \approx 370$ m。

\textbf{8.} $x = \frac{h}{\tan 40^\circ} \approx
\frac{370}{0.839} \approx 441$ m。从 $A$ 处检验：
$\frac{h}{x + 200} \approx \frac{370}{641} \approx 0.577 =
\tan 30^\circ$，二者一致。

\textbf{9.} 单站法需要的是到
\emph{山顶正下方的点}的水平距离 ---
峡谷之间或山体内部根本无法测量。两站法则以一段
你自己能够步量的距离代替它，
即你自己两个观测点之间的 $200$ m。

\textbf{10.} $h = \dfrac{200}{\sqrt3 - \frac{1}{\sqrt3}}
= \dfrac{200}{\frac{2}{\sqrt3}} = 100\sqrt3 \approx 173$ m ---
埃菲尔铁塔的二层，通过两次视测
与一段散步测得。

\textbf{11.} 视线长 $d$，半径 $R$（到切点，
与视线垂直）与 $R + h$（到眼睛）：
由毕达哥拉斯定理得 $d^2 + R^2 = (R + h)^2 = R^2 + 2Rh + h^2$，于是
$d^2 = 2Rh + h^2 \approx 2Rh$（$h$ 与 $R$ 相比极小）。
取眼睛高度 $1.80$ m：
$d \approx \sqrt{2 \times 6.37 \times 10^6 \times 1.8}
\approx 4\,800$ m --- 海平线仅约五
公里远。

\textbf{12.} $d \approx \sqrt{2 \times 6.37 \times 10^6 \times
324} \approx 64$ km。水手的经验法则：$3.6\sqrt{1.8} \approx
4.8$ km 与 $3.6\sqrt{324} = 64.8$ km --- 两者相符。

\textbf{13.} $\tan\theta = \frac{2}{60}$ 得 $\theta \approx
1.9^\circ$。月球的 $0.5^\circ$ 大约是拇指宽度的
\emph{四分之一}：月亮看起来巨大，却能被指甲
挡住 --- 眼睛是一个糟糕的量角器。

\textbf{14.} 楼梯：$\tan^{-1}\frac{17}{29} \approx 30^\circ$
--- 舒适的楼梯倾角几乎恰为
特殊值中的 $30^\circ$。若倾角取 $60^\circ$，踏面宽 $29$ cm，
则踢面高为 $29\tan 60^\circ = 29\sqrt3 \approx
50$ cm：这是梯子，而非楼梯。

\textbf{15.} $\tan 10^\circ \approx 0.18$;
$\tan 45^\circ = 1$; $\tan 80^\circ \approx 5.7$;
$\tan 89^\circ \approx 57.3$。当 $\theta \to 90^\circ$ 时
邻边收缩为零，对边保持不变：
商无限增大，超过任何界限。在 $90^\circ$ 时恰好
不存在三角形，也没有值 --- 正切函数的图像
沿一条竖直渐近线上升，这是
高中分册中众多渐近线的第一条。
\end{solution}
